\documentclass[a4paper,11pt]{article}
\usepackage{fancyhdr,graphics,graphicx,xy}
\usepackage{blkarray}
\usepackage{ifthen}
\usepackage{amssymb,amsmath,mathtools,color}
\usepackage{amsthm}
\pagestyle{fancy}
\usepackage{xparse}


%\textwidth = 6.5 in
%\textheight = 9 in
%\oddsidemargin = 0.0 in
%\evensidemargin = 0.0 in
%\topmargin = 0.0 in
%\headheight = 0.0 in
%\headsep = 0.0 in
%\parskip = 0.2in
%\parindent = 0.0in


%%%%%%%%%%%%%%%%%%%%%%%%%%%%%%%%%%%%%%%%%%
%%% General math macros
%%%%%%%%%%%%%%%%%%%%%%%%%%%%%%%%%%%%%%%%%%

\newcommand{\tf}{\tfrac} 
\renewcommand{\Re}{\mathrm{Re}}

\newcommand{\Def}{\coloneqq}
\newcommand{\lgg}{\operatorname{lg}}
\newcommand{\R}{\mathbb{R}}
\newcommand{\C}{\mathbb{C}}
\newcommand{\D}{\mathbb{D}}
\newcommand{\A}{\mathcal{A}}
\newcommand{\F}{\mathcal{F}}
\newcommand{\Q}{\mathbb{Q}}
\newcommand{\N}{\mathbb{N}} 
\newcommand{\Quat}{\mathbb{J}}
\newcommand{\Z}{\mathbb{Z}}
\newcommand{\T}{\mathbb{T}}
\newcommand{\U}{\mathcal{U}}
\newcommand{\LP}{\operatorname{L}}
\newcommand{\abs}[1]{\left| {#1} \right|}
\newcommand{\sgn}{\operatorname{sgn}}

\newcommand{\tendsto}{\to}
\newcommand{\tendto}{\to}

%%%%%%%%%%%%%%%%%%%%%%%%%%%%%%%%%%%%%%%%%%
%%% Probability Macros
%%%%%%%%%%%%%%%%%%%%%%%%%%%%%%%%%%%%%%%%%%
\newcommand{\stdcg}{N_\bbc(0,1)}

\newcommand{\expec}{\mathbb{E}}
\newcommand{\expect}{\mathbb{E}}
\newcommand{\Exp}{\mathbb{E}}
\newcommand{\E}{\mathbb{E}}
\newcommand{\prob}{\mathbb{P}}
\newcommand{\filt}{\mathscr{F}}
\newcommand{\Prob}{\prob}
\renewcommand{\Pr}{\prob}
\newcommand{\Ent}{\operatorname{Ent}}
\newcommand{\Var}{\operatorname{Var}}
\newcommand{\Cov}{\operatorname{Cov}}
\newcommand{\dtv}{d_{TV}}
\newcommand{\weakto}{\Rightarrow}
\newcommand{\lawof}{\mathscr{L}}
\newcommand{\Bernoulli}{\operatorname{Bernoulli}}
\newcommand{\Binomial}{\operatorname{Binom}}
\newcommand{\Poisson}{\operatorname{Poisson}}
\newcommand{\Normal}{\operatorname{Normal}}
\newcommand{\Unif}{\operatorname{Unif}}
\newcommand{\As}{\ensuremath{\text{a.s.}}}
\newcommand{\distequals}{\overset{\mathrm{d}}{=}}


\newcommand{\asstnumber}{3}
\newcommand{\assigntop}{
  \fancyhead[L]{}
  \fancyhead[C]{\Large Math 550 -- Winter 2024\\ Assignment \asstnumber}
  \fancyhead[R]{ }
}

\newtheorem{theorem}{Theorem}

\newtheorem{corollary}[theorem]{Corollary}
\newtheorem{claim}[theorem]{Claim}
\newtheorem{definition}{Definition}
\newcommand{\p}[1]{\mathbb{P}(#1)}
\newcommand{\px}[1]{\mathbb{P}_x(#1)}
\newcommand{\py}[1]{\mathbb{P}_y(#1)}
\newcommand{\I}[1]{\mathbf{1}_{#1}}
\newcommand{\e}{\mathbb{E}}

\begin{document}
\assigntop


\section*{Question 1}

Intuition is that each iteration of forming $S(X)$ "fills in" the space between points of $X$, and that after a certain number if iterations, there are no gaps left in any direction in the space, resulting in a convex set. Intuition leads us to try to show that $S_{\lceil \log_2 d+1 \rceil}(X) = \text{conv}(X)$. \\
\begin{claim}$S_k$ corresponds to convex combinations of at most $2^{k}$ points of $X$. 	
\end{claim}
The proof follows immediate from the claim since any $x \in \text{conv}(X)$ is a convex combination of at most $d+1$ points of $X$ (Carath\'{e}odory), thus letting $k = \lceil \log_2(d+1) \rceil$, $S_k$ covers all convex combination of at most $2^{\log_2(d+1)} = d+1 $ points of $X$, thus $\text{conv}(X) \subseteq S_{\lceil \log_2 (d+1) \rceil} (X)$. Conversely, $S_k(X)$ for any $k$ contains only convex combinations of points in $X$, thus $ S_{\lceil \log_2 (d+1) \rceil} (X) \subseteq \text{conv}(X) $. Since $S(X)$ contains convex combinations of $2$ points of $X$, $S_2(X)$ contains convex combinations of $4$ points of $X$, and so on. 
\begin{proof}[Proof of claim]
	Base case is simple, consider $k=1$, clear that $S(X)$ contains all convex combinations of $2$ points of $X$, because $S(X)$ contains all line segments between pairs of points in $X$. The induction step is as follows, suppose $x$ is a convex combination of at most $2^k$ points of $X$, then $\exists \lambda_1, \ldots, \lambda_{2^k}$ such that $\displaystyle \sum_{i=1}^{2^k} \lambda_i =1$ and $\lambda_i \geq 1$ $\forall i$, and 
	 \begin{align*}
	 	x &= \lambda_1 x_1 + \lambda_2 x_2 + \ldots + \lambda_{2^k} x_{2^k} \\
			&= (\lambda_1 x_1 + \ldots + \lambda_{2^{k-1}} x_{2^{k-1}}) + (\lambda_{2^{k-1}+1} x_{2^{k-1}+1} + \ldots \lambda_{2^k} x_{2^k})
	 \end{align*}
	 letting $\displaystyle A = \sum_{i=1}^{2^{k-1}} \lambda_i$ and $\displaystyle B = \sum_{i=2^{k-1}+1}^{2^k} \lambda_i$ 
	 then 
	 $$x =A \cdot \overbrace{(\frac{\lambda_1}{A} x_1 + \ldots + \frac{\lambda_{2^{k-1}}}{A} x_{2^{k-1}})}^{\in S_{k-1}(X)} + B\cdot \overbrace{(\frac{\lambda_{2^{k-1}+1}}{B} x_{2^{k-1}+1} + \ldots \frac{\lambda_{2^k}}{B} x_{2^k})}^{\in S_{k-1}(X)}$$
	 where $A + B = 1$, , thus $x \in S_k(X)$. 
	 
	
\end{proof}
\section*{Question 2}
\subsection*{(a)}
$x,y,z \in \mathbb{R}^2$ form a triangle, where the vertices have pairwise distance at most 1. Any triangle such triangle can be embedded inside a equilateral triangle, call it $\Delta$ with all vertices having pairwise distance 1 (shrink the edges of $\Delta$ accordingly). Therefore, it's enough to show  it for $\Delta$. Taking the circumcenter of $\Delta$, call it $m$, the unique point equidistant to all three vertices. Since it's found by taking the point at which all perpendicular bisectors of $\Delta$ (60 degree angles) intersect, then the distance from $m$ to any other vertex is calculated as such 
\begin{align*}
	\sqrt{\left(\frac{\tan 30}{2}\right)^2 + \left(\frac{1}{2}\right)^2} &= \sqrt{ \left(\frac{\sqrt{3}}{3\cdot 2}\right)^2 + \left(\frac{1}{2}\right)^2}\\
	&= \sqrt{\frac{3}{36} + \frac{1}{4}}\\
	&= \sqrt{\frac{1}{3}}\\
	&= \frac{1}{\sqrt{3}}
\end{align*}
Thus $\Delta \subseteq \overline{B(m,\frac{1}{\sqrt{3}})}$ the closed disc centered at $m$ with radius $\frac{1}{\sqrt{3}}$. 
\subsection*{(b)}
Let $C_x$ be the disc centered at $x$ of radius $\frac{1}{\sqrt{3}}$. We know that for any $\{x,y,z\} \subseteq X$, $C_x \cap C_y \cap C_z \neq \emptyset$. Indeed, by (a), for any three pts $\{x,y,z\} \subseteq X$, there exists some point $m$  such that $\|m-x\|, \|m-y\|, \|m-z\| \leq \frac{1}{\sqrt{3}}$, thus $m \in C_x \cap C_y \cap C_z$. By Helly's, since $\{x,y,z\}$ is arbitrary, then $\displaystyle \bigcap_{x \in X} C_x \neq \emptyset$, thus there exists some $m \in \mathbb{R}^2$ such that $\|m - x\| \leq \frac{1}{\sqrt{3}}$ for any $x \in X$, and the result follows. 
\subsection*{(c)}
We first do the case with $4$ points, where any $4$ points in $\mathbb{R}^3$ form a hyperplane in $\mathbb{R}^2$. Consider the extremal case first, which is the regular tetrahedron with edge length 1. It has a circumcircle of radius $\sqrt{\frac{8}{3}}$, thus using the the sam argument as (b), every finite set with diameter at most $1$ can be contained in a closed ball of radius $\sqrt{\frac{8}{3}}$. Then the goal is to show that $\forall \varepsilon > 0$, then finding a counter example for $c = \sqrt{\frac{8}{3}}-\varepsilon $ which I don't know how to do, I've never taken a course in discrete geometry before.

\section*{3}



\end{document}